% !PRINTED 353
\weberpage{353}
% [sic] Die Quelle führt zunächst \\mathfrak{m} ein, spricht aber vor seiner Definition von \\mathfrak{a}; der Wortlaut bleibt erhalten.
%<<CONT

2. Satz. \gesperrt{Jede Zahl der Ordnung $[Q]$ und nur diese sind in der Form enthalten}
\weq{6}{\omega = x + y\theta,}
\gesperrt{worin $x, y$ beliebige ganze rationale Zahlen sind.}

Wir nennen das System
\weq{7}{(1, \theta)}
eine Basis der Ordnung $[Q]$.

Durch Einsetzen von (5) kann man die Zahlen (6) auch in der Form darstellen:
\weq{8}{\omega = \dfrac{x + y\sqrt{D}}{2},}
wobei die $x, y$ ganze Zahlen sind, die der Bedingung genügen
\weq{9}{x^2 - y^2 D \equiv 0 \pmod{4}.}
Die Zahl $D$ ist eine Diskriminante, deren Stamm $\Delta$ ist. Wir nennen sie die \gesperrt{Diskriminante der Ordnung $[Q]$}.

\begin{center}
\S\ 97. \gesperrt{Ordnungen und Ideale.}
\end{center}

Es sei $\mathfrak{m}$ ein Ideal des quadratischen Körpers $\Omega$. Wenn alle durch $\mathfrak{a}$ teilbaren ganzen Zahlen durch eine rationale Zahl $m$ teilbar sind, so ist das Ideal $\mathfrak{a} = m\mathfrak{m}$ durch $m$ teilbar; es ist
\[ N(\mathfrak{a}) = m^2 N(\mathfrak{m}), \]
und die rationale Zahl $m N(\mathfrak{m})$, die kleiner ist als $N(\mathfrak{a})$, ist durch $\mathfrak{a}$ teilbar.

Ist $m$ die größte natürliche Zahl, durch die $\mathfrak{a}$ teilbar ist, so wollen wir $m$ \gesperrt{den Teiler von $\mathfrak{a}$} nennen. Ist der Teiler $= 1$, so heißt das Ideal \gesperrt{primär}, sonst \gesperrt{abgeleitet}. Man erhält alle aus einem primären Ideal $\mathfrak{m}$ abgeleiteten Ideale, indem man $\mathfrak{m}$ mit beliebigen natürlichen Zahlen multipliziert, und alle aus einem primären Ideal abgeleiteten Ideale sind untereinander \gesperrt{äquivalent}.

3. Satz. \gesperrt{Ist $\mathfrak{a}$ ein primäres Ideal, so ist die Norm von $\mathfrak{a}$ zugleich die kleinste durch $\mathfrak{a}$ teilbare natürliche Zahl.}

Bezeichnen wir nämlich die kleinste durch $\mathfrak{a}$ teilbare natürliche Zahl mit $a$, so ist die Norm von $\mathfrak{a}$ jedenfalls durch $a$ teilbar. Wir setzen
\[ N(\mathfrak{a}) = ma, \]
und wenn $\mathfrak{a}'$ das zu $\mathfrak{a}$ konjugierte Ideal ist, so ergibt sich, da nach Bd.\ II, \S\ 164 die Norm eines Ideals gleich dem Produkt der konjugierten Ideale ist,
\weq{1}{N(\mathfrak{a}) = \mathfrak{a}\mathfrak{a}' = ma.}
